\subsection*{10.5 Application: Continuous Mapping Theorem}

In this section we prove two continuous mapping theorems--one for almost sure convergence and one for convergence in probability. The two theorems say that almost sure convergence and convergence in probability are preserved by continuous functions.

\begin{thmbox}{10.11}
Let $(\Omega,\mathcal{F},P)$ be a probability space, let $V$ be a $d$-dimensional random vector, and let $(V_n)_{n=1}^{\infty}$ be a sequence of $d$-dimensional random vectors. Suppose $f:\mathbb{R}^d\to\mathbb{R}^m$ is a function that is continuous on a set $S_f\subseteq\mathbb{R}^d$ with $P(V\in S_f)=1$. Then:
\begin{enumerate}
\item $V_n\xrightarrow{\mathrm{a.s.}}V$ implies $f(V_n)\xrightarrow{\mathrm{a.s.}}f(V)$.
\item $V_n\xrightarrow{P}V$ implies $f(V_n)\xrightarrow{P}f(V)$.
\end{enumerate}
\end{thmbox}

\textit{Proof}
We let $f_k$ be the $k$-th component of the function $f$, for $k=1,2,\ldots,m$.

For part (a), by Theorem 10.10, it is sufficient to show that for each component function $f_k$ of $f$, we have
\[
f_k(V_n)\xrightarrow{\mathrm{a.s.}}f_k(V).
\]
From the hypotheses in the theorem, there exists a set $B$ with probability $1$ such that $V_n(\omega)\to V(\omega)$ for all $\omega\in B$. Let $A$ be the pre-image $V^{-1}(S_f)$. The set $A$ has probability $1$ by assumption. Hence, in $A\cap B$, we have
\[
\lim_{n\to\infty}V_n(\omega)=V(\omega),
\]
and
\[
\lim_{n\to\infty}f_k(V_n(\omega))=f_k(V(\omega)),
\]
by the continuity of $f_k$. Therefore $f_k(V_n)$ converges to $f_k(V)$ for all $\omega$ in $A\cap B$, which has probability $1$.

For part (b), by Theorem 10.10, it is sufficient to show that for each component function $f_k$ of $f$, we have $f_k(V_n)$ converging to $f_k(V)$ in probability. In the following we fix an index $k$.

Let $\delta$ and $\epsilon$ be positive real numbers. Consider the set
\[
A_{\delta,\epsilon}\coloneqq
\{v\in\mathbb{R}^d:\exists w\in\mathbb{R}^d
\text{ such that }\lVert w-v\rVert<\delta
\text{ and }
\lVert f_k(w)-f_k(v)\rVert>\epsilon\}.
\]
For each vector $v$ that is not in $A_{\delta,\epsilon}$, we cannot find any vector $w$ that satisfies $\lVert w-v\rVert<\delta$ and $\lVert f_k(w)-f_k(v)\rVert>\epsilon$ simultaneously. In set notation, we can express it as
\[
\{\omega:\lVert f_k(V_n(\omega))-f_k(V(\omega))\rVert>\epsilon\}
\subseteq
\{\omega:V(\omega)\in A_{\delta,\epsilon}\}
\cup
\{\omega:\lVert V_n(\omega)-V(\omega)\rVert\geq\delta\}.
\]
Take $\delta=1/j$ for positive integer $j$. By the continuity of $f$, and hence of the component function $f_k$, the set $A_{1/j,\epsilon}\cap S_f$ decreases to $\emptyset$ as $j$ tends to infinity. By the upper semi-continuous property of measure, we obtain
\[
P(\{\omega:V(\omega)\in A_{1/j,\epsilon}\})
=
P(\{\omega:V(\omega)\in A_{1/j,\epsilon}\cap S_f\})
\to 0
\]
as $j\to\infty$.

For an arbitrarily small $\epsilon_1>0$, we can find a sufficiently large $M$ such that
\[
P(\{\omega:V(\omega)\in A_{1/\ell,\epsilon}\})\leq \epsilon_1/2
\]
for all $\ell\geq M$.

Since it is assumed that $V_n\to V$ in probability, we can find a sufficiently large $N$ such that
\[
P(\{\lVert V_\ell(\omega)-V(\omega)\rVert\geq\delta\})\leq \epsilon_1/2
\]
for all $\ell\geq N$. This yields
\[
P(\{\lVert f_k(V_\ell(\omega))-f_k(V(\omega))\rVert>\epsilon\})
\leq \epsilon_1/2+\epsilon_1/2=\epsilon_1
\]
for all $\ell\geq \max(M,N)$. As $\epsilon_1$ is arbitrary, the probability
\[
P(\lVert f_k(V_\ell)-f_k(V)\rVert>\epsilon)
\]
converges to $0$ as $\ell\to\infty$.
\hfill $\square$

\begin{thmbox}{10.12}
\begin{enumerate}
\item Suppose $X_n\xrightarrow{\mathrm{a.s.}}X$ and $Y_n\xrightarrow{\mathrm{a.s.}}Y$. Then:
\begin{enumerate}
\item $X_n+Y_n\xrightarrow{\mathrm{a.s.}}X+Y$.
\item $X_nY_n\xrightarrow{\mathrm{a.s.}}XY$.
\item $\alpha X_n\xrightarrow{\mathrm{a.s.}}\alpha X$ for any constant $\alpha$.
\item $X_n/Y_n\xrightarrow{\mathrm{a.s.}}X/Y$ provided that $Y_n$ and $Y$ are nonzero with probability $1$.
\end{enumerate}
\item If $X_n\xrightarrow{P}X$ and $Y_n\xrightarrow{P}Y$, then (a) to (d) above hold with ``a.s.'' replaced by ``$P$''.
\end{enumerate}
\end{thmbox}

\textit{Proof}
We prove part (a) for convergence in probability only. The proofs of the others are similar. Since $X_n$ and $Y_n$ are both converging to $X$ and $Y$, respectively, in probability, the random vector $(X_n,Y_n)$ converges to the random vector $(X,Y)$ in probability. We apply Theorem 10.11 with the continuous function $f(x,y)=x+y$ (cf. the proof of Theorem 4.6).
\hfill $\square$

Exercise 10.10 gives a partial extension of this theorem.

\textbf{Example 10.5.1 (Consistent Estimator)} \\
Given $n$ i.i.d. samples $X_1,X_2,\ldots,X_n$ drawn from a probability distribution with fixed but unknown parameter $\theta$, we estimate $\theta$ by a function
\[
\widehat{\theta}_n(X_1,\ldots,X_n)
\]
of the samples. The estimator $\widehat{\theta}_n$ is called strongly consistent if $\widehat{\theta}_n(X_1,\ldots,X_n)$ converges to $\theta$ a.s. and is called weakly consistent if it converges to $\theta$ in probability.

Suppose a consistent estimator, either strongly or weakly, of the variance of the distribution is available. If we want to estimate the standard deviation, a natural solution is to take the square root of the estimated variance. The continuous mapping theorem guarantees that the resulting estimate is consistent.
